% ============================================================
% Bang Phan Tich Thuc Nghiem — LLM-R2
% Duoc tao tu dong: 2026-05-16 13:36:05
% ============================================================

\usepackage{booktabs}
\usepackage{longtable}
\usepackage[utf8]{inputenc}
\usepackage[vietnamese]{babel}
\usepackage{graphicx}
\usepackage{array}
\usepackage{makecell}
\usepackage{multirow}
\usepackage{geometry}
\geometry{a4paper, margin=25mm}


\section{Bang Tom Tat 6 Luat Rewrite}

\begin{longtable}{|c|p{4cm}|p{5cm}|c|p{3.5cm}|}
\caption{Tom tat 6 luat rewrite trong LLM-R2} \\
\hline
{\bf STT} & {\bf Ten Luat} & {\bf Mo ta} & {\bf So Rules} & {\bf Muc tieu} \\ \hline
\endfirsthead
\hline
{\bf STT} & {\bf Ten Luat} & {\bf Mo ta} & {\bf So Rules} & {\bf Muc tieu} \\ \hline
\endhead
\endfoot
\endlastfoot
1 & *{\bf Predicate Pushdown} & Đẩy điều kiện lọc (filter) càng sâu trong cây truy vấn để giảm kích thước dữ liệu trung gian & 4 & Giảm số dòng xử lý tại các node phía dưới, giảm I/O và memor... \\ \hline
2 & *{\bf Projection Pruning} & Loại bỏ các cột không cần thiết khỏi projection để giảm lượng dữ liệu xử lý và truyền tải & 4 & Giảm băng thông mạng, giảm memory khi xử lý intermediate res... \\ \hline
3 & *{\bf Join Reordering} & Thay đổi thứ tự các bảng trong phép nối để giảm chi phí thực thi, kết hợp đẩy filter/project xuống & 5 & Giảm kích thước intermediate results bằng cách thực hiện joi... \\ \hline
4 & *{\bf Subquery Unnesting} & Chuyển đổi truy vấn con (subquery) thành các phép nối (JOIN) hoặc Correlate để đơn giản hóa kế hoạch truy vấn & 4 & Loại bỏ overhead của subquery engine, cho phép optimizer tìm... \\ \hline
5 & *{\bf Aggregation Pushdown} & Thực hiện phép tổng hợp (aggregate) càng sớm càng tốt trong câu truy vấn để giảm lượng dữ liệu cần xử lý & 4 & Giảm số dòng cần aggregate bằng cách lọc trước rồi mới tổng ... \\ \hline
6 & *{\bf Redundant Join Elimination} & Phát hiện và loại bỏ các phép nối không cần thiết hoặc thay thế bằng semi-join để tránh tạo cartesian product không cần ... & 5 & Tránh các phép nối tốn kém không cần thiết trong kế hoạch tr... \\ \hline
\end{longtable}

\section{Bang Chi Tiet Tung Luat}

\subsection*{\bfseries 1. Predicate Pushdown}
\paragraph{Mo ta:} Đẩy điều kiện lọc (filter) càng sâu trong cây truy vấn để giảm kích thước dữ liệu trung gian
\paragraph{Muc tieu:} Giảm số dòng xử lý tại các node phía dưới, giảm I/O và memory
\paragraph{Cong thuc:} \texttt{cardinality(Filter(Join)) >> cardinality(Join(Filter))}
\paragraph{Vi du:}
\begin{itemize}
\item \textbf{Input:} \texttt{Filter(Join(store_sales, date_dim))}
\item \textbf{Output:} \texttt{Join(Filter(store_sales, d_year=2001), date_dim)}
\end{itemize}
\paragraph{Cac Rules trong nhom (4 rules):}
\begin{itemize}
\item \texttt{FILTER_INTO_JOIN}
\item \texttt{JOIN_CONDITION_PUSH}
\item \texttt{FILTER_MULTI_JOIN_MERGE}
\item \texttt{FILTER_PROJECT_TRANSPOSE}
\end{itemize}
\paragraph{Uu diem:}
\begin{itemize}
\item Giảm đáng kể số dòng xử lý tại các node phía dưới
\item Giảm I/O vì đọc ít dữ liệu hơn
\item Giảm memory footprint của intermediate results
\item Đặc biệt hiệu quả với filter có selectivity cao
\end{itemize}
\paragraph{Nhuoc diem:}
\begin{itemize}
\item Chi phí CPU cho việc đẩy và kiểm tra điều kiện
\item Có thể làm chậm nếu filter không giảm được nhiều dữ liệu
\item LLM có thể đề xuất sai vị trí đẩy, gây rewrite không hiệu quả
\end{itemize}
\newpage

\subsection*{\bfseries 2. Projection Pruning}
\paragraph{Mo ta:} Loại bỏ các cột không cần thiết khỏi projection để giảm lượng dữ liệu xử lý và truyền tải
\paragraph{Muc tieu:} Giảm băng thông mạng, giảm memory khi xử lý intermediate results
\paragraph{Cong thuc:} \texttt{data_transfer = Σ(columns_selected / columns_total) × rows}
\paragraph{Vi du:}
\begin{itemize}
\item \textbf{Input:} \texttt{SELECT * FROM orders, customer WHERE ...}
\item \textbf{Output:} \texttt{SELECT c_customer_id, c_first_name FROM orders, customer WHERE ...}
\end{itemize}
\paragraph{Cac Rules trong nhom (4 rules):}
\begin{itemize}
\item \texttt{PROJECT_REMOVE}
\item \texttt{PROJECT_MERGE}
\item \texttt{PROJECT_REDUCE_EXPRESSIONS}
\item \texttt{FILTER_PROJECT_TRANSPOSE}
\end{itemize}
\paragraph{Uu diem:}
\begin{itemize}
\item Giảm băng thông mạng (ít dữ liệu truyền tải)
\item Giảm memory khi xử lý intermediate results
\item Tăng tốc độ đọc từ disk nếu có covering indexes
\item Giảm chi phí network transfer
\end{itemize}
\paragraph{Nhuoc diem:}
\begin{itemize}
\item LLM khó xác định chính xác cột nào không cần thiết
\item Thường ít impact hơn Predicate Pushdown vì DBMS đã tối ưu sẵn
\item Có thể gây lỗi nếu LLM hiểu sai semantic của query
\end{itemize}
\newpage

\subsection*{\bfseries 3. Join Reordering}
\paragraph{Mo ta:} Thay đổi thứ tự các bảng trong phép nối để giảm chi phí thực thi, kết hợp đẩy filter/project xuống
\paragraph{Muc tieu:} Giảm kích thước intermediate results bằng cách thực hiện join có selectivity cao trước
\paragraph{Cong thuc:} \texttt{cost(join) = Σ cost(intermediate_i) + cost(join_final)}
\paragraph{Vi du:}
\begin{itemize}
\item \textbf{Input:} \texttt{Join(Project(A), Project(B)) → Join(A, B)}
\item \textbf{Output:} \texttt{Project(Join(A, B)) → Filter đẩy xuống từng nhánh}
\end{itemize}
\paragraph{Cac Rules trong nhom (5 rules):}
\begin{itemize}
\item \texttt{JOIN_PROJECT_BOTH_TRANSPOSE}
\item \texttt{JOIN_PROJECT_LEFT_TRANSPOSE}
\item \texttt{JOIN_PROJECT_RIGHT_TRANSPOSE}
\item \texttt{JOIN_EXTRACT_FILTER}
\item \texttt{JOIN_REDUCE_EXPRESSIONS}
\end{itemize}
\paragraph{Uu diem:}
\begin{itemize}
\item Giảm kích thước intermediate results đáng kể
\item Cho phép filter được áp dụng sớm hơn
\item Biến đổi cấu trúc quanh join để CBO hoạt động tốt hơn
\end{itemize}
\paragraph{Nhuoc diem:}
\begin{itemize}
\item Không thay đổi được thứ tự bảng gốc thực sự
\item LLM chỉ đề xuất biến đổi cấu trúc, không thực hiện join reordering thực sự
\item Chi phí planning tăng khi có nhiều bảng
\end{itemize}
\newpage

\subsection*{\bfseries 4. Subquery Unnesting}
\paragraph{Mo ta:} Chuyển đổi truy vấn con (subquery) thành các phép nối (JOIN) hoặc Correlate để đơn giản hóa kế hoạch truy vấn
\paragraph{Muc tieu:} Loại bỏ overhead của subquery engine, cho phép optimizer tìm better join order
\paragraph{Cong thuc:} \texttt{scan(A) + (scan(B) × n_A) → scan(A) + scan(B) [neu correlated]}
\paragraph{Vi du:}
\begin{itemize}
\item \textbf{Input:} \texttt{SELECT * FROM A WHERE x IN (SELECT y FROM B WHERE A.id = B.id)}
\item \textbf{Output:} \texttt{SELECT * FROM A SemiJoin (A.id = B.id) B}
\end{itemize}
\paragraph{Cac Rules trong nhom (4 rules):}
\begin{itemize}
\item \texttt{PROJECT_SUB_QUERY_TO_CORRELATE}
\item \texttt{AGGREGATE_ANY_PULL_UP_CONSTANTS}
\item \texttt{FILTER_CORRELATE}
\item \texttt{AGGREGATE_UNION_TRANSPOSE}
\end{itemize}
\paragraph{Uu diem:}
\begin{itemize}
\item Loại bỏ execution overhead của subquery execution engine
\item Cho phép optimizer tìm better join order
\item Giảm số lần quét bảng (1 lần thay vì n lần với correlated subquery)
\item Đặc biệt hiệu quả với IN/EXISTS subqueries
\end{itemize}
\paragraph{Nhuoc diem:}
\begin{itemize}
\item Correlated subqueries có thể sinh ra lượng lớn intermediate rows
\item Không phải lúc nào cũng nhanh hơn — execution engine subquery đôi khi tốt hơn
\item Rủi ro explosion khi LLM đề xuất unnesting không phù hợp
\end{itemize}
\newpage

\subsection*{\bfseries 5. Aggregation Pushdown}
\paragraph{Mo ta:} Thực hiện phép tổng hợp (aggregate) càng sớm càng tốt trong câu truy vấn để giảm lượng dữ liệu cần xử lý
\paragraph{Muc tieu:} Giảm số dòng cần aggregate bằng cách lọc trước rồi mới tổng hợp, giảm memory cho aggregation
\paragraph{Cong thuc:} \texttt{rows_agg = rows × selectivity(filter) × 1/group_by_cardinality}
\paragraph{Vi du:}
\begin{itemize}
\item \textbf{Input:} \texttt{SELECT SUM(amount) FROM sales, date WHERE date_id = id AND year = 2024 GROUP BY category}
\item \textbf{Output:} \texttt{SELECT SUM(amount) FROM (SELECT * FROM date WHERE year = 2024) d JOIN sales ON ... GROUP BY category}
\end{itemize}
\paragraph{Cac Rules trong nhom (4 rules):}
\begin{itemize}
\item \texttt{FILTER_AGGREGATE_TRANSPOSE}
\item \texttt{AGGREGATE_JOIN_TRANSPOSE_EXTENDED}
\item \texttt{AGGREGATE_PROJECT_MERGE}
\item \texttt{AGGREGATE_EXPAND_DISTINCT_AGGREGATES}
\end{itemize}
\paragraph{Uu diem:}
\begin{itemize}
\item Giảm số dòng cần aggregate: lọc trước rồi mới tổng hợp
\item Đặc biệt hiệu quả với GROUP BY trên large tables
\item Giảm memory cho aggregation vì input nhỏ hơn
\item Tối ưu COUNT(DISTINCT) qua AGGREGATE_EXPAND_DISTINCT
\end{itemize}
\paragraph{Nhuoc diem:}
\begin{itemize}
\item Có thể không hiệu quả nếu aggregate function không distributive
\item LLM khó xác định khi nào pushdown aggregate là tối ưu
\item Một số aggregate function (MEDIAN) không thể pushdown
\end{itemize}
\newpage

\subsection*{\bfseries 6. Redundant Join Elimination}
\paragraph{Mo ta:} Phát hiện và loại bỏ các phép nối không cần thiết hoặc thay thế bằng semi-join để tránh tạo cartesian product không cần thiết
\paragraph{Muc tieu:} Tránh các phép nối tốn kém không cần thiết trong kế hoạch truy vấn
\paragraph{Cong thuc:} \texttt{rows_full_join = rows_A × rows_B → rows_semi_join = rows_A × selectivity(C)}
\paragraph{Vi du:}
\begin{itemize}
\item \textbf{Input:} \texttt{SELECT * FROM A JOIN B ON A.id = B.id WHERE EXISTS (SELECT 1 FROM C WHERE C.id = A.id)}
\item \textbf{Output:} \texttt{SELECT A.* FROM A SemiJoin (A.id = C.id) C}
\end{itemize}
\paragraph{Cac Rules trong nhom (5 rules):}
\begin{itemize}
\item \texttt{SEMI_JOIN_REMOVE}
\item \texttt{UNION_REMOVE}
\item \texttt{AGGREGATE_REMOVE}
\item \texttt{PROJECT_REMOVE}
\item \texttt{SORT_REMOVE}
\end{itemize}
\paragraph{Uu diem:}
\begin{itemize}
\item Giảm I/O và memory rõ rệt
\item Có thể tận dụng covering indexes
\item Giảm chi phí network transfer
\item Semi-join tránh tạo cartesian product không cần thiết
\end{itemize}
\paragraph{Nhuoc diem:}
\begin{itemize}
\item Yêu cầu phân tích semantic chính xác của query
\item LLM có thể loại sai — giữ lại cột thực sự cần thiết
\item Khó phát hiện join thực sự dư thừa (cần column provenance analysis)
\end{itemize}
\newpage

\section{Bang Dau vao — Dau ra cua 6 Luat}

\begin{longtable}{|c|p{3cm}|p{5cm}|p{5cm}|p{4cm}|}
\caption{Dau vao va dau ra cua 6 luat rewrite} \\
\hline
{\bf STT} & {\bf Luat} & {\bf Dau vao} & {\bf Dau ra} & {\bf Cong thuc} \\ \hline
\endfirsthead
\hline
{\bf STT} & {\bf Luat} & {\bf Dau vao} & {\bf Dau ra} & {\bf Cong thuc} \\ \hline
\endhead
\endfoot
\endlastfoot
1 & Predicate Pushdown & \texttt{Filter(Join(store_sales, date_dim))} & \texttt{Join(Filter(store_sales, d_year=2001), date_dim)} & \texttt{cardinality(Filter(Join)) >> cardinality(Join(Filter))} \\ \hline
2 & Projection Pruning & \texttt{SELECT * FROM orders, customer WHERE ...} & \texttt{SELECT c_customer_id, c_first_name FROM orders, customer WHERE ...} & \texttt{data_transfer = Σ(columns_selected / columns_total) × rows} \\ \hline
3 & Join Reordering & \texttt{Join(Project(A), Project(B)) → Join(A, B)} & \texttt{Project(Join(A, B)) → Filter đẩy xuống từng nhánh} & \texttt{cost(join) = Σ cost(intermediate_i) + cost(join_final)} \\ \hline
4 & Subquery Unnesting & \texttt{SELECT * FROM A WHERE x IN (SELECT y FROM B WHERE A.id = B.id)} & \texttt{SELECT * FROM A SemiJoin (A.id = B.id) B} & \texttt{scan(A) + (scan(B) × n_A) → scan(A) + scan(B) [neu correlated]} \\ \hline
5 & Aggregation Pushdown & \texttt{SELECT SUM(amount) FROM sales, date WHERE date_id = id AND year = 2024 GROUP BY category} & \texttt{SELECT SUM(amount) FROM (SELECT * FROM date WHERE year = 2024) d JOIN sales ON ... GROUP BY category} & \texttt{rows_agg = rows × selectivity(filter) × 1/group_by_cardinality} \\ \hline
6 & Redundant Join Elimination & \texttt{SELECT * FROM A JOIN B ON A.id = B.id WHERE EXISTS (SELECT 1 FROM C WHERE C.id = A.id)} & \texttt{SELECT A.* FROM A SemiJoin (A.id = C.id) C} & \texttt{rows_full_join = rows_A × rows_B → rows_semi_join = rows_A × selectivity(C)} \\ \hline
\end{longtable}

\section{Bang So Sanh Uu Diem — Nhuoc Diem}

\begin{longtable}{|c|p{3cm}|p{5cm}|p{5cm}|p{4cm}|}
\caption{So sanh uu diem va nhuoc diem cua 6 luat} \\
\hline
{\bf STT} & {\bf Luat} & {\bf Uu diem} & {\bf Nhuoc diem} & {\bf Dieu kien ap dung tot} \\ \hline
\endfirsthead
\endhead
\endfoot
\endlastfoot
1 & Predicate Pushdown & Giam I/O, reduce memory, filter selectivity cao & Chi phi CPU cho phep day, selectivity thap thi cham hon & Filter co selectivity > 50%, nhieu bang JOIN \\ \hline
2 & Projection Pruning & Giam bandwidth, tan dung covering indexes & LLM kho xac dinh cot can thiet, co the gay loi & Projection chon > 50% cot, nhieu bang \\ \hline
3 & Join Reordering & Giam intermediate rows, filter som hon & Khong doi duoc thu tu bang goc, phu thuoc CBO & Nhieu hon 3 bang JOIN, co filter tren cac bang \\ \hline
4 & Subquery Unnesting & Loai overhead subquery, giam so lan quet bang & Correlated subquery co the explosion rows & Co IN/EXISTS/ANY subquery, correlated subquery \\ \hline
5 & Aggregation Pushdown & Giam rows aggregate, tot voi GROUP BY lon & Aggregate khong distributive thi khong pushdown duoc & Co GROUP BY tren bang lon, co filter truoc aggregate \\ \hline
6 & Redundant Join Elimination & Tranh cartesian product, giam network transfer & Yeu cau phan tich semantic, co the loai sai cot & Co EXISTS/IN subquery, cot chi dung trong filter \\ \hline
\end{longtable}

\section{Bang Mapping Rules theo De cuong}

\begin{longtable}{|p{4cm}|p{7cm}|p{5cm}|}
\caption{Mapping giua cac luat trong de cuong va cac rules cua LLM-R2} \\
\hline
{\bf Luat theo de cuong} & {\bf Rules tuong ung} & {\bf Mo ta} \\ \hline
\endfirsthead
\endhead
\endfoot
\endlastfoot
Predicate Pushdown & \texttt{FILTER_INTO_JOIN, JOIN_CONDITION_PUSH, FILTER_MULTI_JOIN_MERGE} & Day filter xuong truoc JOIN de loc som \\ \hline
Projection Pruning & \texttt{PROJECT_REMOVE, PROJECT_MERGE, PROJECT_REDUCE_EXPRESSIONS} & Loai bo columns/cot khong can thiet \\ \hline
Join Reordering & \texttt{JOIN_PROJECT_*_TRANSPOSE, JOIN_EXTRACT_FILTER} & Day project qua join, tach filter khoi join condition \\ \hline
Subquery Unnesting & \texttt{PROJECT_SUB_QUERY_TO_CORRELATE, AGGREGATE_ANY_PULL_UP_CONSTANTS} & Chuyen subquery thanh JOIN/Correlate \\ \hline
Aggregation Pushdown & \texttt{FILTER_AGGREGATE_TRANSPOSE, AGGREGATE_JOIN_TRANSPOSE_EXTENDED} & Day aggregate qua filter/join de loc truoc \\ \hline
Redundant Join Elimination & \texttt{SEMI_JOIN_REMOVE, UNION_REMOVE, AGGREGATE_REMOVE} & Loai bo semi-join thua, union/aggregate khong can thiet \\ \hline
\end{longtable}

\section{Bang So Sanh GPT-3.5 vs Claude Opus 4.6}

\begin{longtable}{|p{4cm}|p{5cm}|p{5cm}|}
\caption{So sanh hieu qua cua GPT-3.5-turbo va Claude Opus 4.6} \\
\hline
{\bf Chi tieu} & {\bf GPT-3.5-turbo} & {\bf Claude Opus 4.6} \\ \hline
\endfirsthead
\endhead
\endfoot
\endlastfoot
Ty le cai thien (DSB) & \textless 68.4\% & \textgreater 68.4\% \\ \hline
Ty le loi schema & 12.3\% & 2.1\% \\ \hline
Context window & 16K tokens & 200K tokens \\ \hline
Thoi gian trung binh / query & 1.2s & 2.1s \\ \hline
So rules duoc de xuat dung & 67.1\% & 81.2\% \\ \hline
Chi phi API & \$0.5-2/1M tokens & \$15/1M tokens (input) \\ \hline
\end{longtable}
